\documentclass[../../main.tex]{subfiles}
% \documentclass[a4paper,10pt]{ltjsarticle}
\begin{document}
\rhead{京都大学後期$1993$年$6$番}
\setcounter{equation}{0}
\begin{oframed}
  \subfile{problem}
\end{oframed}
\setlength{\columnseprule}{0.4pt}
\begin{multicols}{2}
  {\bf [解答]}
  与えられた不等式を通分すると
  \begin{align}
    & \dfrac{a_1x+b_1}{x+c_1} > \dfrac{a_2x+b_2}{x+c_2}   \\
    & \dfrac{(a_1 x + b_1)(x + c_2) - (a_2 x + b_2)(x + c_1)}{(x + c_1)(x + c_2)} > 0 \label{eq:1}
  \end{align}
  である．これが$x\neq -c_1,-c_2$で常に成り立つような条件を考える．
  この分子を $f(x)$ とおくと，
  \begin{align}
    f(x) = (a_1 - a_2)x^2 + (a_1 c_2 + b_1 - b_2 - a_2 c_1)x + b_1 c_2 - b_2 c_1
  \end{align}
  であり，これは$x$の二次以下の式である．
  まず，$a_1-a_2=0$の時は$f(x)$は一次式となり，\cref{eq:1}が成り立たない$x$が必ず存在するから，
  \begin{align}
    a_1-a_2 \neq 0\label{eq:2}
  \end{align}
  であることが必要．
  この元で，\cref{eq:1}は分母分子ともに二次式であり，さらに分母が$x=-c_1,-c_2$を根をもつ．
  よって\cref{eq:1}が$x\neq -c_1,-c_2$の全ての実数$x$で成り立つには，
  分子の二次係数が正であり，かつ分子が分母の定数倍になっていなければならない．
  \begin{align}
    \begin{dcases}
      a_1-a_2 > 0 \\
      f(x) = (a_1-a_2)(x+c_1)(x+c_2)
    \end{dcases}
  \end{align}
  $f(x)$の一次と零次の係数比較すると
  \begin{align}
    \begin{dcases}
      a_1 c_2 + b_1 - b_2 - a_2 c_1 & = (a_1 - a_2)(c_1 + c_2) \\
      b_1 c_2 - b_2 c_1             & = (a_1 - a_2)c_1 c_2
    \end{dcases}
  \end{align}
  整理すると，
  \begin{align}
    &
    \begin{dcases}
      a_1 c_1 - a_2 c_2 = b_1 - b_2 \\
      (a_1 - a_2)c_1 c_2 = b_1 c_2 - b_2 c_1
    \end{dcases}
    &
    \begin{dcases}
      b_1 - a_1 c_1 = b_2 - a_2 c_2 \\
      c_1(b_2 - a_2 c_2) = c_2(b_1 - a_1 c_1)
    \end{dcases} \label{eq:3}
  \end{align}
  となる．
  \cref{eq:3}を$c_1 = c_2$ か否かで場合分けすると，

  \textbf{1$^\circ$ $c_1 = c_2$の時}
  \begin{align}
    b_2 - a_2 c_2 = b_1 - a_1 c_1
  \end{align}
  が求める条件となる．

  \textbf{2$^\circ$ $c_1 \ne c_2$の時}
  \begin{align}
    b_2 - a_2 c_2 = b_1 - a_1 c_1 = 0
  \end{align}
  が求める条件となる．

  以上から，題意が成立する条件は\cref{eq:2}と合わせて，
  \begin{align}
    c_1 = c_2 \land b_2 - a_2 c_2 = b_1 - a_1 c_1 \land a_1 > a_2 \label{eq:4}
  \end{align}
  又は
  \begin{align}
    c_1 \ne c_2 \land b_2 = a_2 c_2 \land b_1 = a_1 c_1 \land a_1 > a_2 \label{eq:5}
  \end{align}
  である．以下それぞれが十分かどうかを確認する．

  \subparagraph{\cref{eq:4}の時}

  \cref{eq:4}より$c_1 = c_2$だから，\cref{eq:1}に代入して
  \begin{align}
    & \dfrac{(a_1-a_2)x + (b_1-b_2)}{x+c_1} > 0
  \end{align}
  である．さらに\cref{eq:4}より$b_1-b_2=a_1c_1-a_2c_2=(a_1-a_2)c_1$だから
  \begin{align}
    & \dfrac{(a_1-a_2)x + (a_1-a_2)c_1}{x+c_1} > 0 \\
    & \dfrac{(a_1-a_2)(x + c_1)}{x+c_1} > 0 \\
    \therefore
    & a_1-a_2 > 0
  \end{align}
  となるが，これは\cref{eq:4}から成立する．よってこの場合は十分である．

  \subparagraph{\cref{eq:5}の時}

  \cref{eq:1}に$b_2 = a_2 c_2, b_1 = a_1 c_1$を代入して
  \begin{align}
    & \dfrac{(a_1 x + a_1c_1)(x + c_2) - (a_2 x + a_2c_2)(x + c_1)}{(x + c_1)(x + c_2)} > 0 \\
    & \dfrac{a_1(x + c_1)(x + c_2) - a_2(x + c_2)(x + c_1)}{(x + c_1)(x + c_2)} > 0  \\
    & (a_1-a_2)>0
  \end{align}
  となるが，これは\cref{eq:5}から成立する．よってこの場合は十分である．
  \casesend

  以上から，題意が成立する条件は
  \begin{align}
    c_1 = c_2 \land b_2 - a_2 c_2 = b_1 - a_1 c_1 \land a_1 > a_2
  \end{align}
  又は
  \begin{align}
    c_1 \ne c_2 \land b_2 = a_2 c_2 \land b_1 = a_1 c_1 \land a_1 > a_2
  \end{align}
  である．  $\cdots$(答)

  {\bf [解説]}

  \newpage
\end{multicols}
\end{document}
